% Источники: «6 Семинар-2.jpg», «6 Семинар-4.jpg», «6 Семинар-5.jpg».
% Условия: conditions/sem_06.pdf. В записях разобраны задачи 1, 4, 5(a).
\section{Функция распределения: смеси распределений}

\noindent Условия:
\href{run:conditions/sem_06.pdf}{семинар 6 (PDF)}.

\begingroup
\everymath{\textstyle}
\subsection{Задача 1}

\begin{solution}
    \textbf{Модель: дискретная.}
    \[
        \Pr\{\xi=k\}=\frac{C}{k(k+1)},\qquad k=1,2,\ldots
    \]
    \textbf{а)} Сумма вероятностей равна $1$:
    \[
        1=C\sum_{k=1}^{\infty}\frac1{k(k+1)}.
    \]
    Разложим слагаемое и найдём частичную сумму:
    \[
        \frac1{k(k+1)}=\frac1k-\frac1{k+1},\qquad
        \sum_{k=1}^{n}\left(\frac1k-\frac1{k+1}\right)
        =1-\frac1{n+1}\longrightarrow1.
    \]
    Следовательно, $\boxed{C=1}$.

    \medskip
    \textbf{б)}
    \[
        \Pr\{\xi\leq3\}
        =\sum_{k=1}^{3}\Pr\{\xi=k\}
        =\frac12+\frac16+\frac1{12}=\boxed{\frac34}.
    \]
\end{solution}

\subsection{Задача 4}

\begin{solution}
    Пусть $F_k$ --- функция распределения $\xi_k$,
    $p_k=\Pr\{\nu=k\}$, $k\in\mathbb N$.
    Случайный индекс $\nu$ не зависит от $\xi_1,\xi_2,\ldots$.
    Величина $\xi_\nu$ определяется поточечно:
    \[
        \nu(\omega)=k\quad\Longrightarrow\quad
        \xi_\nu(\omega)=\xi_k(\omega).
    \]
    По формуле полной вероятности и независимости:
    \[
        \begin{aligned}
            F_{\xi_\nu}(x)
            &=\Pr\{\xi_\nu\leq x\}\\
            &=\sum_{k:\,p_k>0}
                \Pr\{\xi_\nu\leq x\mid\nu=k\}\,p_k\\
            &=\sum_{k:\,p_k>0}
                \Pr\{\xi_k\leq x\mid\nu=k\}\,p_k\\
            &=\sum_{k=1}^{\infty}F_k(x)\,p_k.
        \end{aligned}
    \]
    Таким образом,
    \[
        \boxed{F_{\xi_\nu}(x)=\sum_{k=1}^{\infty}p_kF_k(x).}
    \]
\end{solution}

\clearpage
\subsection{Разложение распределения}

\begin{remark}
    Любую функцию распределения можно представить в виде
    \[
        F(x)=\alpha_1F_{\text{д}}(x)
            +\alpha_2F_{\text{а.н.}}(x)
            +\alpha_3F_{\text{с}}(x),
        \qquad
        \alpha_i\geq0,\quad \alpha_1+\alpha_2+\alpha_3=1,
    \]
    где $F_{\text{д}}$ --- дискретная,
    $F_{\text{а.н.}}$ --- абсолютно непрерывная,
    $F_{\text{с}}$ --- сингулярная непрерывная функция распределения.
    Слагаемые с нулевыми коэффициентами опускаются.
\end{remark}

\subsection{Задача 5(a)}

\begin{solution}
    Функция распределения, заданная графиком:
    \[
        F(x)=\begin{cases}
            0,&x<0,\\
            1/4,&0\leq x<1,\\
            1/2,&1\leq x<2,\\
            (x-1)/2,&2\leq x<3,\\
            1,&x\geq3.
        \end{cases}
    \]
    \begin{center}
        \begin{tikzpicture}[x=1.7cm,y=3cm,>=Stealth]
            \draw[->] (-0.65,0)--(3.7,0) node[right] {$x$};
            \draw[->] (0,-0.12)--(0,1.22) node[above] {$F(x)$};
            \foreach \x in {1,2,3}
                \draw (\x,0.02)--(\x,-0.02) node[below] {$\x$};
            \foreach \y/\lab in {0.25/{1/4},0.5/{1/2},1/{1}}
                \node[left] at (0,\y) {$\lab$};
            \node[below left] at (0,0) {$0$};
            \draw[dashed] (0,1)--(3,1)--(3,0);
            \draw[blue,very thick] (-0.6,0)--(0,0);
            \draw[blue,very thick] (0,0.25)--(1,0.25);
            \draw[blue,very thick] (1,0.5)--(2,0.5)--(3,1)--(3.6,1);
            \foreach \x/\y in {0/0,1/0.25}
                \filldraw[blue,fill=solution_back] (\x,\y) circle (1.7pt);
            \foreach \x/\y in {0/0.25,1/0.5,2/0.5,3/1}
                \fill[blue] (\x,\y) circle (1.7pt);
        \end{tikzpicture}
    \end{center}
    Скачки в точках $0$ и $1$ равны $1/4$ каждый.
    Суммарная масса дискретной части равна $1/2$;
    после нормировки получаем $\operatorname{Bern}(1/2)$.

    На отрезке $[2,3]$ функция линейно возрастает от $1/2$ до $1$.
    Масса абсолютно непрерывной части также равна $1/2$;
    после нормировки получаем $U[2,3]$.
    \[
        \boxed{F(x)=\frac12 F_{\operatorname{Bern}(1/2)}(x)
                    +\frac12 F_{U[2,3]}(x).}
    \]
\end{solution}
\endgroup
